\section*{第 \ref{ch:g6:light-propagation} 章——光的直线传播}

\begin{solution}{exo:g6:light-propagation:1}
在清澈的\omterm{def:g2:air-around-us:air}{空气}中（或任何清澈、混合均匀的介质中），光
沿直线传播。“直线传播”中，“直线”即沿直线，“传播”即
光从一处到另一处的行进。
\end{solution}

\begin{solution}{exo:g6:light-propagation:2}
发散：\omterm{def:g6:light-propagation:ray}{光线}从一点向外散开（如蜡烛、裸露的\omterm{def:g3:first-electric-circuit:bulb}{灯泡}）。
平行：\omterm{def:g6:light-propagation:ray}{光线}并排前进（阳光——太阳极其遥远，
其\omterm{def:g6:light-propagation:ray}{光线}到达我们时几乎并排行进）。会聚：\omterm{def:g6:light-propagation:ray}{光线}向一点
靠拢（如放大镜汇聚的光）。
\end{solution}

\begin{solution}{exo:g6:light-propagation:3}
\omterm{def:g6:light-propagation:ray}{光线}可以当作直尺：眼睛、边缘与远端
三者恰好对齐，即远端的光恰好
擦着整条边缘到达眼睛——这依据的正是\omterm{prop:g6:light-propagation:law}{光的直线传播}。
\end{solution}

\begin{solution}{exo:g6:light-propagation:4}
恰好擦过挡光物边缘的\omterm{def:g6:light-propagation:ray}{光线}，其落点
即为\omterm{def:g1:light-and-shadows:shadow}{影子}的边界。位于\omterm{def:g1:light-and-shadows:shadow}{影子}锥体之内的眼睛看不到灯——
因为从灯发出的光中，没有一束能沿直线到达它。
\end{solution}

\begin{solution}{exo:g6:light-propagation:5}
位于中间时，\omterm{def:g1:light-and-shadows:shadow}{影子}为小球的两倍，即 \qty{12}{cm}；靠近
屏幕时，\omterm{def:g1:light-and-shadows:shadow}{影子}缩小，趋向小球的真实
大小 \qty{6}{cm}。
\end{solution}

\begin{solution}{exo:g6:light-propagation:6}
屏幕底部处：自高处一点穿过
低处小门的直线，必然继续向下延伸。
\end{solution}

\begin{solution}{exo:g6:light-propagation:7}
上下颠倒、左右互换：屏幕上，手的挥动
显示在低处而非高处，且位于相反的一侧。两种颠倒
都源自\omterm{def:g6:light-propagation:ray}{光线}在小孔处交叉。
\end{solution}

\begin{solution}{exo:g6:light-propagation:8}
天文学家：可在屏幕上安全观察日食，而不必
直视天空。画家：描摹投影中的景物，以掌握透视。
（现代的同类：每一台照相机，乃至眼睛本身。）
\end{solution}

\begin{solution}{exo:g6:light-propagation:9}
\omterm{def:g1:light-and-shadows:shadow}{影子}缩小，趋向硬币的真实大小：灯
离得更远时，擦过边缘的\omterm{def:g6:light-propagation:ray}{光线}离开的角度更平缓——
扇形更窄，落点也更靠近。作图
表明：\omterm{def:g6:light-propagation:ray}{光线}更平缓，展开更紧凑。
\end{solution}

\begin{solution}{exo:g6:light-propagation:10}
太阳小而遥远：近似平行的\omterm{def:g6:light-propagation:ray}{光线}构成紧凑的扇形，
边界清晰——\omterm{def:g1:light-and-shadows:shadow}{影子}鲜明。灰色云层则是一盏
与整个天空一样大的灯：光从各个方向同时射来，
地面每一点都有部分\omterm{def:g6:light-propagation:ray}{光线}到达，\omterm{def:g1:light-and-shadows:shadow}{影子}
几乎消散殆尽。
\end{solution}

\begin{solution}{exo:g6:light-propagation:11}
更亮：宽大的小孔让每个景物点射入的\omterm{def:g6:light-propagation:ray}{光线}更多。
模糊：经过宽孔时，每个景物点的\omterm{def:g6:light-propagation:ray}{光线}会到达屏幕上
整块的\emph{区域}，而非一个点，无数
重叠的区域使图像变得模糊。针孔的清晰
正在于它的\emph{小}。
\end{solution}

\begin{solution}{exo:g6:light-propagation:12}
树叶之间的每一道缝隙都是天然的针孔；地面上每一块
亮斑都是该针孔所成的太阳的像——普通日子里是圆的，
因为太阳本身是圆的。日食时太阳变成月牙形，
每道树叶缝隙都会忠实地投出
月牙形的像：数百台小照相机，拍摄的是同一个天体。
\end{solution}

\begin{solution}{pb:g6:light-propagation:1}
\textbf{1.} 紧贴屏幕时，擦过边缘的\omterm{def:g6:light-propagation:ray}{光线}已无距离
可以展开：\omterm{def:g1:light-and-shadows:shadow}{影子}紧贴木偶自身的轮廓。
\textbf{2.} 位于中间时\omterm{def:g1:light-and-shadows:shadow}{影子}加倍：\qty{40}{cm}。
\textbf{3.} 四倍的 $20$ 恰为 \qty{80}{cm}，与
之相符；木偶位于灯与屏幕之间
距离的四分之一处。
\textbf{4.} 大而发光的球体相当于同时有许多盏灯：每一部分
都投出各自略有错位的\omterm{def:g1:light-and-shadows:shadow}{影子}，龙的
边缘变得柔和而灰暗。清晰的影戏需要一个小而裸露、
近似点状的\omterm{def:g1:light-and-shadows:source}{光源}。
\textbf{5.} 靠近灯时：擦过边缘的\omterm{def:g6:light-propagation:ray}{光线}所成的扇形变宽，
狼影变大；手臂低垂，不在扇形之内。木偶
靠近那盏小而实在的灯时，边缘柔化略有加剧——
近距离时，灯的宽度的影响更大。
\textbf{6.} 从灯出发穿过两个木偶的直线，必须
落在屏幕上不同的区域——两个木偶左右相隔
足够远，使得任一木偶都不在另一个木偶的
影锥之内。
\textbf{7.} 手的\omterm{def:g1:light-and-shadows:shadow}{影子}：手离灯更近，位于全程的
三分之一处，其擦过边缘的\omterm{def:g6:light-propagation:ray}{光线}所成扇形为三倍宽，而
位于中段的龙只有两倍——靠近灯时，即使是小手也能
扮演巨人。
\textbf{8.} 把反派紧贴屏幕放置，它的\omterm{def:g1:light-and-shadows:shadow}{影子}
缩小为本身的大小（或完全移出\omterm{def:g6:light-propagation:ray}{光束}）；或把
第二个挡光物靠近灯——它巨大的\omterm{def:g1:light-and-shadows:shadow}{影子}会吞没
反派的\omterm{def:g1:light-and-shadows:shadow}{影子}，反派便“消失”在另一片
黑暗之中。
\textbf{9.} 在纸的低处：从高处的灯出发，穿过孔洞的
\omterm{def:g6:light-propagation:ray}{光线}会继续向下延伸——针孔
使像倒置。
\textbf{10.} 从右到左：小孔使左右
与上下一同互换。
\textbf{11.} 更亮，却模糊得无可救药：街道上的每一点
如今都投出硬币大小的光斑，光斑相互重叠，
画面消散——大结局随清晰一同消逝。
\textbf{12.}“今晚的全部节目——每一道\omterm{def:g1:light-and-shadows:shadow}{影子}、每一个巨人，
以及一条上下颠倒的街道——都由同一条定律上演：
光沿直线传播。”
\end{solution}
